\section{RF and M/W Network Theory and Analysis}
{\color{sub}\footnotesize 4 hrs \gap$\cdot$\gap asked in \mk{18 of 24 papers}
\gap$\cdot$\gap 2 syllabus sub-topics \gap$\cdot$\gap \mk{193 marks, 10.1\% of the archive}
$\rightarrow$ definitions, one derivation, four matrix properties and one device
(the magic tee) that answers half the chapter.}

\vspace{3pt}\noindent
\colorbox{gdL}{\makebox[\dimexpr\textwidth-2\fboxsep][l]{%
  \textcolor{gd}{\bfseries\footnotesize READ THIS FIRST}}}
\par\nobreak\vspace{3pt}

Chapter 3 is the cheapest 8 to 10 marks in the paper, because the same three asks recur:
\begin{enumerate}[label=\arabic*., leftmargin=1.4em, itemsep=1pt]
  \item \textbf{Why S-parameters, and how do they differ from h / Z / Y / ABCD?}
    (13 papers) $\rightarrow$ one reason list plus one comparison table, $\S$3.1.
  \item \textbf{Derive the S-parameters of a two-port network} (7 papers)
    $\rightarrow$ the same six lines every time, $\S$3.2. Learn it as a script.
  \item \textbf{Identify the passive device from this S-matrix} (5 papers)
    $\rightarrow$ \mk{in every one of the five it is the magic tee}, in four different
    disguises. Learn the fingerprint in $\S$3.7 and the matrix in $\S$3.6.
\end{enumerate}
Two more traps worth knowing before the paper starts:
\begin{itemize}
  \item The radar questions (``two transmitters, twice the power'', ``half the power to the
    antenna via a duplexer'', 4 papers) are the \textbf{same magic tee} asked as a design
    problem, $\S$3.8. Nothing new has to be learnt for them.
  \item An examiner who writes ``3-port network'' almost always wants the
    \textbf{impossibility theorem} ($\S$3.5), not just a $3\times3$ matrix. Stating that a
    3-port cannot be lossless, reciprocal and matched at the same time, and then naming the
    three escapes, is what separates a 4 from an 8.
\end{itemize}

\hr

%=======================================================================
\T{3.1 Why S-parameters: the Limits of Z, Y, h and ABCD}

\Q{Why are S-parameters used in microwave network analysis? Compare S-parameters with h-parameters}{\m{2}\m{3}\m{4}\m{4+4}\m{5+5}}
{\color{sub}\footnotesize \tS{13} \yr{69 Bh, \textbf{70 Bh}, 70 Ma, 73 Ma, \textbf{73 Bh}, 74 Ma, \textbf{75 Bh}, \textbf{78 Ch}, \textbf{\texttt{79 Bh}}, \texttt{80 Ba}, \textbf{\texttt{80 Bh}}, \textbf{80 Ch}, \textbf{\texttt{81 Bh}}} \gap$\cdot$\gap \mk{the single most repeated ask in the chapter}}

\sbs{%
\lead{What the low-frequency parameter sets need}
\begin{itemize}
  \item All four classical sets relate \textbf{total voltage and total current} at the
    terminals:
  \begin{itemize}
    \item $[V] = [Z][I]$, $Z_{ij} = V_i/I_j$ with all other ports \textbf{open}, $I_k = 0$.
    \item $[I] = [Y][V]$, $Y_{ij} = I_i/V_j$ with all other ports \textbf{shorted}, $V_k = 0$.
    \item h-parameters: mixed, $h_{11}$ and $h_{21}$ need a \textbf{short} on port 2,
      $h_{12}$ and $h_{22}$ need an \textbf{open} on port 1.
    \item ABCD: cascade form, still defined on $V$ and $I$.
  \end{itemize}
  \item Two things are therefore assumed:
  \begin{itemize}
    \item $V$ and $I$ are \textbf{unique} at a terminal pair (true only when
      $d \ll \lambda$).
    \item Ideal \textbf{open} and \textbf{short} terminations are available.
  \end{itemize}
  \item Both assumptions fail above roughly $1\,$GHz.
\end{itemize}}{%
\lead{What S-parameters do instead}
\begin{itemize}
  \item Defined on \textbf{travelling waves}, not on totals:
  \begin{itemize}
    \item $a_i$ = wave incident on port $i$, $b_i$ = wave leaving port $i$.
    \item Normalised so $|a_i|^2$ and $|b_i|^2$ \textbf{are powers}, in watts.
  \end{itemize}
  \item Every port is terminated in $Z_0$, a \textbf{matched load}:
  \begin{itemize}
    \item Easy to build and broadband, unlike open or short.
    \item No reflection back into the device $\rightarrow$ no oscillation, no damage,
      bias undisturbed.
  \end{itemize}
  \item Measured with \textbf{directional couplers} plus a vector network analyser, which
    separate incident from reflected wave and read magnitude and phase directly.
  \item Every practical figure of merit falls straight out: gain, return loss, insertion
    loss, isolation, VSWR, stability ($\S$3.4).
  \item Valid at \textbf{any} frequency, so the same formalism covers DC to mmWave.
  \item Cascadable after conversion to the transfer (T) form, so a chain of measured
    blocks can be multiplied out.
\end{itemize}}

\sbs{%
\lead{The five reasons they break at microwave}
\begin{enumerate}[label=\arabic*., leftmargin=1.4em]
  \item \textbf{No unique $V$ and $I$}:
  \begin{itemize}
    \item Circuit size is comparable to $\lambda$, so voltage varies along the conductor.
    \item In a waveguide there is no unique pair of conductors at all, so ``the voltage''
      has no single definition.
  \end{itemize}
  \item \textbf{Open and short cannot be realised}:
  \begin{itemize}
    \item A physical short has series lead inductance, an open has fringing capacitance.
    \item A short is a short at \textbf{one frequency and one reference plane} only:
      $\lambda/4$ away it looks like an open.
  \end{itemize}
  \item \textbf{Total reflection is destructive}:
  \begin{itemize}
    \item $|\Gamma| = 1$ sends all the power back into the device.
    \item Transistors and negative-resistance diodes \textbf{oscillate or burn out} under
      open or short test conditions.
  \end{itemize}
  \item \textbf{No instrument measures RF $V$ and $I$}:
  \begin{itemize}
    \item Probes load the circuit and cannot follow GHz waveforms.
    \item \textbf{Power} is what a microwave detector actually reads.
  \end{itemize}
  \item \textbf{Bias shifts}:
  \begin{itemize}
    \item Open and short terminations change the DC operating point and the junction
      capacitance of active devices, so the measurement is not of the device in use.
  \end{itemize}
\end{enumerate}}{%
\lead{h-parameters vs S-parameters}
\begin{tabularx}{\linewidth}{@{}L{2.1cm}XX@{}}
\toprule
\hrow \thead{Aspect} & \thead{h-parameters} & \thead{S-parameters} \\
\midrule
\textbf{Variables}   & Total $V$, $I$ & Incident, reflected waves $a$, $b$ \\
\textbf{Termination} & Short and open & Matched $Z_0$ \\
\textbf{Measured}    & Voltage, current & Power (magnitude + phase) \\
\textbf{Instrument}  & Curve tracer, bridge & Vector network analyser \\
\textbf{On active devices} & Oscillation, damage & Safe, bias preserved \\
\textbf{Frequency}   & Below $\approx100$ MHz & Any, chosen for RF / M/W \\
\textbf{Broadband}   & No, termination drifts & Yes \\
\textbf{Cascading}   & Awkward & Via T-matrix \\
\bottomrule
\end{tabularx}}

\penalty0\vspace{2pt}

\creamq{Answer skeleton for the 4-mark version}{\m{4}}
\begin{itemize}
  \item One line on what fails: no unique $V$ / $I$, open and short unrealisable, active
    devices damaged, no RF voltmeter.
  \item One line on the fix: normalised travelling waves measured against matched $Z_0$
    loads with directional couplers.
  \item One line on the payoff: gain, RL, IL, VSWR and stability read straight off $[S]$.
  \item Close with $[b] = [S][a]$ and the two-port matrix. Four sentences, full marks.
\end{itemize}

\penalty0\vspace{2pt}

%=======================================================================
\T{3.2 Definition and Derivation for a Two-Port Network}

\Q{Using a two-port network, define and derive the S-parameters}{\m{4}\m{5}\m{6}\m{10}}
{\color{sub}\footnotesize \tS{7} \yr{69 Bh, \textbf{70 Bh}, 70 Ma, \textbf{78 Ch}, \texttt{80 Ba}, \textbf{\texttt{80 Bh}}, \textbf{80 Ch}} \gap$\cdot$\gap \mk{write it as a script: figure, normalisation, matrix, four conditions}}

\sbs{%
\lead{Step 1: the wave variables}
\begin{itemize}
  \item At port $i$ the line carries a forward wave $V_i^{+}$ and a backward wave
    $V_i^{-}$. Normalise both by $\sqrt{Z_0}$:
    \[ a_i = \frac{V_i^{+}}{\sqrt{Z_0}}, \qquad b_i = \frac{V_i^{-}}{\sqrt{Z_0}} \]
  \item The normalisation is chosen so that the squared magnitudes are powers:
    \[ P_i^{\text{inc}} = \tfrac{1}{2}|a_i|^2, \qquad
       P_i^{\text{ref}} = \tfrac{1}{2}|b_i|^2 \]
    (the $\tfrac{1}{2}$ disappears if $a$, $b$ are taken as rms).
  \item In terms of the total quantities:
    \[ V_i = \sqrt{Z_0}\,(a_i + b_i), \qquad I_i = \frac{a_i - b_i}{\sqrt{Z_0}} \]
    so nothing is lost, only re-expressed.
\end{itemize}
\figC{c3_twoport.png}{0.85\linewidth}
{\footnotesize\color{sub} Port 1 and port 2 of a general two-port, each carrying an
incident wave $a$ and an emerging wave $b$. Source: RF Pulchowk Ch3 deck.}}{%
\lead{Step 2: the defining matrix}
\begin{itemize}
  \item The network is linear, so each outgoing wave is a weighted sum of all the incoming
    ones:
    \[ b_1 = S_{11}a_1 + S_{12}a_2, \qquad b_2 = S_{21}a_1 + S_{22}a_2 \]
    \[ \begin{bmatrix} b_1 \\ b_2 \end{bmatrix} =
       \begin{bmatrix} S_{11} & S_{12} \\ S_{21} & S_{22} \end{bmatrix}
       \begin{bmatrix} a_1 \\ a_2 \end{bmatrix}
       \quad\text{i.e.}\quad [b] = [S][a] \]
  \item \textbf{Subscript order}: in $S_{ij}$ the \textbf{second} index $j$ is the port
    driven, the \textbf{first} index $i$ is the port observed. $S_{21}$ is
    ``out of 2, in at 1''.
  \item Written out, the matrix is the complete input-output description of the black box:
    knowing four numbers replaces solving for every internal voltage and current.
\end{itemize}}

\sbs{%
\lead{Step 3: isolate each parameter}
\begin{itemize}
  \item Terminate the unused port in $Z_0$. A matched load reflects nothing, so its
    incident wave vanishes and one column of the matrix survives at a time.
\end{itemize}
\begin{tabularx}{\linewidth}{@{}L{2.4cm}L{1.5cm}X@{}}
\toprule
\hrow \thead{Parameter} & \thead{Condition} & \thead{Meaning} \\
\midrule
$S_{11} = \dfrac{b_1}{a_1}$ & $a_2 = 0$ & Input reflection coefficient, port 2 matched \\
$S_{21} = \dfrac{b_2}{a_1}$ & $a_2 = 0$ & Forward transmission (gain or insertion loss) \\
$S_{12} = \dfrac{b_1}{a_2}$ & $a_1 = 0$ & Reverse transmission (isolation, feedback) \\
$S_{22} = \dfrac{b_2}{a_2}$ & $a_1 = 0$ & Output reflection coefficient, port 1 matched \\
\bottomrule
\end{tabularx}
\vspace{2pt}
\begin{itemize}
  \item \mk{$a_2 = 0$ means ``port 2 terminated in $Z_0$'', never ``port 2 open''.}
    Writing ``open circuit'' here loses the mark, because it is exactly the condition
    S-parameters were invented to avoid.
\end{itemize}
\vspace{2pt}
\figC{c3_sfg.png}{0.60\linewidth}
{\footnotesize\color{sub} Signal-flow graph of the same two-port: four branches, one per
S-parameter. Useful when the question asks for a ``flow diagram''.}}{%
\lead{Step 4: tie $S_{11}$ back to impedance}
\begin{itemize}
  \item With port 2 matched, port 1 simply sees the input impedance $Z_{in}$, so
    \[ S_{11} = \Gamma_{in} = \frac{Z_{in} - Z_0}{Z_{in} + Z_0} \]
  \item This is the bridge to Chapter 2: a one-port load has the single entry
    $[S] = [\,\Gamma_L\,]$, and a perfectly matched two-port has
    $[S] = \left[\begin{smallmatrix}0&1\\1&0\end{smallmatrix}\right]$ ($\S$2.8).
\end{itemize}
\lead{Step 5: generalise to $n$ ports}
\begin{itemize}
  \item $b_i = \sum_{j=1}^{n} S_{ij}a_j$, an $n \times n$ matrix with $n^2$ entries:
    \[ S_{ij} = \left.\frac{b_i}{a_j}\right|_{a_k = 0,\ k \neq j} \]
  \item Rows and columns both equal the port count, so a 3-port gives $3\times3$ and a
    magic tee gives $4\times4$.
  \item \textbf{Reference impedance}: every entry is defined against the same $Z_0$
    (usually $50\,\Omega$), so a quoted matrix without its $Z_0$ and its frequency is
    incomplete.
\end{itemize}}

\penalty0\vspace{2pt}

%=======================================================================
\T{3.3 Properties of the S-Matrix}

\Q{State and justify the properties of the scattering matrix}{\m{4}\m{4+4}\m{5}\m{6}}
{\color{sub}\footnotesize \tS{8} \yr{\textbf{73 Bh}, \textbf{74 Bh}, \textbf{75 Bh}, \textbf{78 Ch}, \textbf{\texttt{79 Bh}}, \textbf{\texttt{80 Bh}}, \textbf{\texttt{81 Bh}}, \textbf{\texttt{82 Bh}}} \gap$\cdot$\gap \mk{these four tests are also the tool that identifies a device from its matrix}}

\sbs{%
\lead{P1. Zero diagonal $\Leftrightarrow$ matched port}
\begin{itemize}
  \item $S_{ii} = 0$ means a wave arriving at port $i$ is not reflected, so that port is
    matched to $Z_0$.
  \item ``Matched at all ports'' means the whole diagonal is zero.
\end{itemize}
\lead{P2. Symmetry $\Leftrightarrow$ reciprocity}
\begin{itemize}
  \item $S_{ij} = S_{ji}$ for all $i, j$, i.e. $[S] = [S]^{T}$.
  \item \textbf{Holds} for networks built of passive, linear, isotropic media: lines,
    stubs, tees, couplers, resistors, capacitors, transformers.
  \item \textbf{Fails} for magnetised ferrite (circulator, isolator, gyrator) and for
    active devices (amplifier, transistor), where $S_{21} \neq S_{12}$ by design.
  \item So a matrix that is \textbf{not} symmetric can never be a passive isotropic device.
\end{itemize}}{%
\lead{P3. Unitary $\Leftrightarrow$ lossless}
\begin{itemize}
  \item A lossless network redirects power but never dissipates it, which forces
    \[ [S]^{*T}[S] = [I] \]
  \item Reading that out column by column gives the two tests actually used in exams:
  \begin{itemize}
    \item \textbf{Column-sum test} (each column carries unit power):
      \[ \sum_{k=1}^{n} |S_{ki}|^2 = 1 \quad \text{for every } i \]
    \item \textbf{Orthogonality test} (distinct columns are orthogonal):
      \[ \sum_{k=1}^{n} S_{ki}\,S_{kj}^{*} = 0 \quad \text{for every } i \neq j \]
  \end{itemize}
  \item If a column sums to less than 1, the shortfall is the fraction of incident power
    the network \textbf{dissipates}, which is how ``is it lossless?'' is answered
    numerically ($\S$3.4 and Problem 1 of the companion).
\end{itemize}}

\sbs{%
\lead{P4. Shifting the reference plane multiplies by a phase}
\begin{itemize}
  \item Move the reference plane of port $i$ outward by an electrical length
    $\theta_i = \beta \ell_i$. Then
    \[ S'_{ij} = S_{ij}\,e^{-j(\theta_i + \theta_j)} \]
  \item Diagonal terms pick up $e^{-2j\theta_i}$, because the wave travels the added
    length twice.
  \item \textbf{Magnitudes are unchanged}: $|S'_{ij}| = |S_{ij}|$. Return loss, insertion
    loss and VSWR do not depend on where the plane is drawn; only phase does.
  \item This is why an S-parameter measurement is meaningless until the planes are
    declared, and why a VNA must be calibrated to a stated plane.
\end{itemize}
\lead{P5. Complex, frequency dependent, $n^2$ entries}
\begin{itemize}
  \item Each $S_{ij}$ carries a magnitude and an angle, and both move with frequency, so
    a quoted set is meaningless without its $f$ and its $Z_0$.
  \item An $n$-port has $n^2$ coefficients; symmetry and structural zeros cut the number
    of independent ones sharply (the magic tee has 16 entries but only 2 distinct values).
\end{itemize}
\lead{P6. Symmetric and antimetric two-ports}
\begin{itemize}
  \item \textbf{Symmetric}: reciprocal \emph{and} $S_{11} = S_{22}$. The network looks the
    same from either end.
  \item \textbf{Antimetric}: $S_{11} = -S_{22}$.
\end{itemize}}{%
\lead{The four tests as an exam checklist}
\begin{tabularx}{\linewidth}{@{}L{2.5cm}L{2.9cm}X@{}}
\toprule
\hrow \thead{Test} & \thead{Look for} & \thead{Conclusion} \\
\midrule
Diagonal & $S_{ii} = 0$? & That port is matched \\
Transpose & $[S] = [S]^{T}$? & Reciprocal (passive) or not (ferrite / active) \\
Column sums & $\sum_k|S_{ki}|^2 = 1$? & Lossless, else lossy by the shortfall \\
Column dots & $\sum_k S_{ki}S_{kj}^{*} = 0$? & Confirms lossless, ties the phases \\
\bottomrule
\end{tabularx}
\vspace{2pt}
{\footnotesize\color{sub} Run the four in this order on any matrix the examiner prints:
they are enough to name the device ($\S$3.7) and to answer every ``judge the condition''
rider.}}

\penalty0\vspace{2pt}

%=======================================================================
\T{3.4 Return Loss, Insertion Loss and the Other Figures of Merit}

\Q{Define return loss and insertion loss, and the losses associated with a two-port network}{\m{2}\m{1}\m{4}}
{\color{sub}\footnotesize \tF{3} \yr{\textbf{72 Ash}, \textbf{\texttt{79 Bh}}, \textbf{\texttt{82 Bh}}} \gap$\cdot$\gap small marks, but they are the language every other answer is written in}

\sbs{%
\lead{The three powers at a port}
\begin{itemize}
  \item $P_i$ incident, $P_r$ reflected, $P_t$ transmitted; whatever is left over is
    dissipated.
  \item Everything below is a ratio of two of those, expressed in dB and rewritten in
    terms of $[S]$.
\end{itemize}
\begin{tabularx}{\linewidth}{@{}L{2.45cm}L{3.1cm}X@{}}
\toprule
\hrow \thead{Quantity} & \thead{In terms of $[S]$} & \thead{What it tells you} \\
\midrule
\textbf{Return loss} & $RL = -20\log_{10}|S_{11}|$ &
  How well port 1 is matched. Bigger is better \\
\textbf{Insertion loss} & $IL = -20\log_{10}|S_{21}|$ &
  Power lost passing through. $0$ dB is ideal \\
\textbf{Isolation} & $-20\log_{10}|S_{12}|$ &
  Leakage backwards. Bigger is better \\
\textbf{Forward gain} & $G = 20\log_{10}|S_{21}|$ &
  Same number as $IL$, opposite sign, used for amplifiers \\
\textbf{VSWR} & $\dfrac{1+|S_{11}|}{1-|S_{11}|}$ &
  The same match on a $1$ to $\infty$ scale \\
\textbf{Reflection (mismatch) loss} & $-10\log_{10}\!\left(1-|S_{11}|^2\right)$ &
  The part of $P_i$ that never enters the network \\
\textbf{Dissipated fraction} & $1 - \sum_k |S_{k1}|^2$ &
  What the network eats. Zero for a lossless network \\
\bottomrule
\end{tabularx}
\vspace{2pt}
\begin{itemize}
  \item \textbf{Useful anchors}: $|S_{11}| = 0.5 \rightarrow RL = 6$ dB, VSWR $=3$;
    $|S_{11}| = 0.1 \rightarrow RL = 20$ dB, VSWR $= 1.22$; $|S_{11}| = 0 \rightarrow$
    perfect match, $RL = \infty$.
\end{itemize}}{%
\figT{c3_sparam_view.png}
{\footnotesize\color{sub} What each S-parameter of a device under test is called in the
lab. Source: RF Pulchowk Ch3 deck.}
\lead{Sign conventions that cost marks}
\begin{itemize}
  \item Return loss and insertion loss are quoted as \textbf{positive} dB for a passive
    device. The minus sign is already in the definition.
  \item $20\log$ for a \textbf{ratio of waves} ($S$ is an amplitude ratio),
    $10\log$ for a \textbf{ratio of powers}. Mixing the two is the commonest slip.
\end{itemize}}

\sbs{%
\lead{Reading a device off its numbers}
\begin{itemize}
  \item Small $|S_{11}|$, $|S_{22}|$ $\rightarrow$ well matched at both ends.
  \item $|S_{21}| > 1$ $\rightarrow$ \textbf{active}: the device has gain.
  \item $|S_{21}| < 1$ with $S_{12} = S_{21}$ $\rightarrow$ \textbf{passive and
    reciprocal}: an attenuator, filter or line.
  \item $|S_{12}| \ll |S_{21}|$ $\rightarrow$ \textbf{unilateral}: an amplifier or
    isolator; the smaller $|S_{12}|$ is, the more stable the amplifier tends to be.
  \item This is the whole method for \textbf{\texttt{82 Bh}}, which prints two amplifier
    S-parameter sets and asks only for ``the basic differences''. Problem 2 of the
    companion works it.
\end{itemize}}{%
\lead{\added{} Correction to the source notes}
\begin{itemize}
  \item The RF Pulchowk deck prints
    $IL\,\text{(dB)} = 10\log\left(|a_1|^2/|b_1|^2\right)$.
  \item $|b_1|$ is the wave coming \textbf{back out of port 1}, so that expression is the
    return loss, not the insertion loss.
  \item Insertion loss compares the incident wave with what leaves the \textbf{far} port:
    \[ IL = 10\log\frac{|a_1|^2}{|b_2|^2} = 20\log\frac{1}{|S_{21}|} \]
    which is what the deck's own next line prints. Use the $|S_{21}|$ form.
\end{itemize}}

\penalty0\vspace{2pt}

%=======================================================================
\T{3.5 Three-Port Networks and Why the Ideal One Cannot Exist}

\Q{Analyse a three-port network using S-parameters. Write the properties of a 3-port network and define its characteristic parameters}{\m{4}\m{4+4}\m{2+6}\m{5+2}}
{\color{sub}\footnotesize \tF{6} \yr{\textbf{72 Ash}, 73 Ma, \textbf{73 Bh}, 74 Ma, \textbf{\texttt{80 Bh}}, \textbf{\texttt{82 Bh}}} \gap$\cdot$\gap \mk{the theorem is the answer, the matrix alone is half marks}}

\sbs{%
\lead{The general 3-port}
\begin{itemize}
  \item Nine coefficients:
    \[ [S] = \begin{bmatrix} S_{11} & S_{12} & S_{13} \\
                             S_{21} & S_{22} & S_{23} \\
                             S_{31} & S_{32} & S_{33} \end{bmatrix} \]
  \item \textbf{Characteristic parameters} an examiner expects you to define:
  \begin{itemize}
    \item \textbf{Matching}: $S_{ii}$, hence return loss and VSWR at each port.
    \item \textbf{Coupling / division ratio}: $|S_{ij}|^2$, the fraction of the input
      power delivered to port $i$.
    \item \textbf{Isolation}: $-20\log|S_{ij}|$ for the pair meant to be decoupled.
    \item \textbf{Phase relation}: $\arg S_{13}$ against $\arg S_{23}$, i.e. whether the
      two output arms are in phase or in antiphase.
    \item \textbf{Insertion loss} and \textbf{bandwidth} over which all of the above hold.
  \end{itemize}
\end{itemize}}{%
\lead{The three escapes, and the device each one gives}
\begin{tabularx}{\linewidth}{@{}L{2.35cm}L{2.5cm}X@{}}
\toprule
\hrow \thead{Give up} & \thead{Device} & \thead{Consequence} \\
\midrule
\textbf{Lossless} & Resistive divider, matched 3-port & All ports matched and reciprocal, but power is burnt in the resistors \\
\textbf{Reciprocal} & \textbf{Circulator} (ferrite) & Matched and lossless, but $S_{ij} \neq S_{ji}$; power circulates $1\!\to\!2\!\to\!3\!\to\!1$ \\
\textbf{Matched at every port} & \textbf{E-plane tee, H-plane tee}, T-junction & Lossless and reciprocal, but only the side arm can be matched \\
\bottomrule
\end{tabularx}
\vspace{2pt}
\begin{itemize}
  \item \textbf{Ideal circulator}, matched and lossless but not reciprocal:
    \[ [S] = \begin{bmatrix} 0&0&1 \\ 1&0&0 \\ 0&1&0 \end{bmatrix} \]
    Terminate port 3 in a matched load and it becomes an \textbf{isolator}.
  \item A \textbf{4-port} escapes the trap entirely: a directional coupler and a magic
    tee are matched, lossless \emph{and} reciprocal all at once, which is exactly why the
    exam favourite is a 4-port ($\S$3.6).
\end{itemize}}

\sbs{%
\lead{Theorem: no 3-port is lossless, reciprocal and matched at once}
\begin{itemize}
  \item \textbf{Assume all three.} Matched gives $S_{11}=S_{22}=S_{33}=0$, reciprocal
    gives $S_{ij}=S_{ji}$, so
    \[ [S] = \begin{bmatrix} 0 & S_{12} & S_{13} \\
                             S_{12} & 0 & S_{23} \\
                             S_{13} & S_{23} & 0 \end{bmatrix} \]
  \item \textbf{Apply the orthogonality test} to each pair of columns:
    \[ S_{13}^{*}S_{23} = 0, \qquad S_{12}^{*}S_{23} = 0, \qquad
       S_{12}^{*}S_{13} = 0 \]
    so at least \textbf{two} of $S_{12}, S_{13}, S_{23}$ must be zero.
  \item \textbf{Apply the column-sum test}:
    \[ |S_{12}|^2 + |S_{13}|^2 = 1, \quad |S_{12}|^2 + |S_{23}|^2 = 1, \quad
       |S_{13}|^2 + |S_{23}|^2 = 1 \]
    which requires all three to be non-zero.
  \item The two demands \textbf{contradict}, so no such network exists.
    \mk{Give up exactly one of the three properties.}
\end{itemize}}{%
\lead{E-plane tee: the worked 3-port ($\S$3.6 has the full matrix)}
\begin{itemize}
  \item Ports 1 and 2 collinear, port 3 the E-arm, assumed matched to the junction.
  \item Three physical facts start the derivation:
  \begin{itemize}
    \item An input at the E-arm splits into two outputs that are $180^\circ$
      \textbf{out of phase}: $S_{23} = -S_{13}$.
    \item Port 3 is matched: $S_{33} = 0$.
    \item The junction is passive and isotropic: $S_{ij} = S_{ji}$.
  \end{itemize}
  \item Four unknowns are left, and the unitary property fixes all four
    ($\S$3.6, and Problem 5 of the companion works every line).
  \item \mk{Ports 1 and 2 cannot be matched}, and indeed the answer gives
    $S_{11} = S_{22} = 1/2$, which is the theorem showing up in the arithmetic.
\end{itemize}}

\penalty0\vspace{2pt}

%=======================================================================
\T{3.6 S-Matrix Catalogue of the Standard Devices}

\Q{Prepare the S-matrices of a perfectly working E-plane Tee, H-plane Tee and a Duplexer, and explore their uses}{\m{4+4}\m{8}\m{3+7}}
{\color{sub}\footnotesize \tF{4} \yr{\textbf{78 Ch}, \texttt{81 Ba}, \textbf{\texttt{80 Bh}}, \textbf{\texttt{82 Bh}}} \gap$\cdot$\gap memorise these five matrices and half the chapter is already answered}

\sbs{%
\lead{E-plane tee (series tee)}
\figT{c3_etee.png}
\begin{itemize}
  \item Side arm in the plane of the \textbf{E} field, so it excites the collinear arms in
    \textbf{antiphase}.
    \[ [S] = \begin{bmatrix}
       \tfrac{1}{2} & \tfrac{1}{2} & \tfrac{1}{\sqrt{2}} \\[2pt]
       \tfrac{1}{2} & \tfrac{1}{2} & -\tfrac{1}{\sqrt{2}} \\[2pt]
       \tfrac{1}{\sqrt{2}} & -\tfrac{1}{\sqrt{2}} & 0 \end{bmatrix} \]
  \item \textbf{Fingerprint}: the $-1/\sqrt{2}$ in the third row / column.
  \item \textbf{Use}: power splitting with a phase reversal; the difference arm of a magic
    tee.
  \item \mk{Ports 1 and 2 are not matched} ($S_{11}=S_{22}=1/2$), which is the 3-port
    theorem of $\S$3.5 showing up in the answer.
\end{itemize}}{%
\lead{H-plane tee (shunt tee)}
\figT{c3_htee.png}
\begin{itemize}
  \item Side arm in the plane of the \textbf{H} field, so the collinear arms come out
    \textbf{in phase}.
    \[ [S] = \begin{bmatrix}
       \tfrac{1}{2} & -\tfrac{1}{2} & \tfrac{1}{\sqrt{2}} \\[2pt]
       -\tfrac{1}{2} & \tfrac{1}{2} & \tfrac{1}{\sqrt{2}} \\[2pt]
       \tfrac{1}{\sqrt{2}} & \tfrac{1}{\sqrt{2}} & 0 \end{bmatrix} \]
  \item \textbf{Fingerprint}: the minus sign has moved to $S_{12}$; the side-arm column is
    all positive.
  \item \textbf{Use}: in-phase power divider or combiner; the sum arm of a magic tee.
  \item \textbf{E-tee vs H-tee in one line}: same geometry, same magnitudes, opposite
    phase relation. The E-tee subtracts, the H-tee adds.
\end{itemize}}

\sbs{%
\lead{Magic tee (hybrid tee): the exam favourite}
\figT{c3_magictee.png}
\begin{itemize}
  \item Ports 1 and 2 collinear, port 3 the \textbf{E-arm} (difference), port 4 the
    \textbf{H-arm} (sum). Four physical facts:
  \begin{itemize}
    \item E-arm drive gives antiphase outputs: $S_{23} = -S_{13}$.
    \item H-arm drive gives in-phase outputs: $S_{24} = +S_{14}$.
    \item E-arm and H-arm are decoupled: $S_{34} = S_{43} = 0$.
    \item Ports 3 and 4 matched: $S_{33} = S_{44} = 0$; and the collinear arms are
      decoupled, $S_{12} = S_{21} = 0$.
  \end{itemize}
  \item The unitary property then forces $S_{11} = S_{22} = 0$ as well, so the tee ends up
    matched at \textbf{all four} ports:
    \[ [S] = \frac{1}{\sqrt{2}}\begin{bmatrix}
       0 & 0 & 1 & 1 \\ 0 & 0 & -1 & 1 \\
       1 & -1 & 0 & 0 \\ 1 & 1 & 0 & 0 \end{bmatrix} \]
  \item \textbf{Fingerprint}: an all-zero $2\times2$ block on each diagonal, one column
    with a sign flip (E-arm), one without (H-arm).
\end{itemize}}{%
\lead{Circulator, isolator, gyrator}
\begin{itemize}
  \item Circulator: $[S]$ as in $\S$3.5, matched and lossless but \textbf{not} symmetric.
  \item Isolator: a circulator with port 3 loaded,
    $[S] = \left[\begin{smallmatrix}0&0\\1&0\end{smallmatrix}\right]$; forward passes,
    reverse is absorbed.
  \item Gyrator: $[S] = \left[\begin{smallmatrix}0&-1\\1&0\end{smallmatrix}\right]$, a
    $180^\circ$ phase difference between the two directions.
\end{itemize}
\lead{Directional coupler}
\begin{itemize}
  \item Through arm $\alpha$, coupled arm $j\beta$, isolated arm 0, with
    $\alpha^2 + \beta^2 = 1$:
    \[ [S] = \begin{bmatrix}
       0 & \alpha & j\beta & 0 \\ \alpha & 0 & 0 & j\beta \\
       j\beta & 0 & 0 & \alpha \\ 0 & j\beta & \alpha & 0 \end{bmatrix} \]
  \item A $3$ dB coupler has $\alpha = \beta = 1/\sqrt{2}$, and is then a magic tee in
    quadrature form.
\end{itemize}}

\sbs{%
\lead{Duplexer}
\begin{itemize}
  \item \textbf{Circulator duplexer} (the usual answer): port 1 transmitter, port 2
    antenna, port 3 receiver. Ideally lossless, so \textbf{all} the transmitter power
    reaches the antenna and \textbf{all} the echo reaches the receiver, with the
    transmitter isolated from the receiver.
  \item \textbf{Magic-tee duplexer}: transmitter on the H-arm, antenna and a matched
    dummy load on the collinear arms, receiver on the E-arm. Splits the transmitter power
    in \textbf{half} and keeps transmitter and receiver isolated. This is the version
    \textbf{76 Bh} asks for ($\S$3.8).
\end{itemize}}{%
\lead{Amplifier, attenuator, amplitude modulator}
\begin{itemize}
  \item Matched amplifier: $[S] = \left[\begin{smallmatrix}0&S_{12}\\S_{21}&0\end{smallmatrix}\right]$
    with $|S_{21}| > 1 \gg |S_{12}|$. Not reciprocal, not lossless.
  \item Matched attenuator of $L$ dB:
    $[S] = \left[\begin{smallmatrix}0&\alpha\\ \alpha&0\end{smallmatrix}\right]$,
    $\alpha = 10^{-L/20}$. Reciprocal but lossy.
  \item \textbf{Amplitude modulator} (\textbf{78 Ch}): a matched, reciprocal two-port
    whose transmission is \textbf{driven by the modulating signal},
    \[ [S] = \begin{bmatrix} 0 & \alpha(t) \\ \alpha(t) & 0 \end{bmatrix},
       \qquad \alpha(t) = \alpha_0\left(1 + m\cos\omega_m t\right) \]
    Problem 6 of the companion derives it and evaluates the numbers.
\end{itemize}}

\penalty0\vspace{2pt}

\creamq{\added{} Magic tee with both E- and H-arms short-circuited}{\m{4+4}}
{\color{sub}\footnotesize \tP{1} \yr{\textbf{\texttt{82 Bh}}} \gap$\cdot$\gap not in any source note, derived here and checked in \code{sparam.py}}

\sbs{%
\begin{itemize}
  \item Shorts at the junction reference plane give $a_3 = -b_3$ and $a_4 = -b_4$.
  \item Substituting into the magic-tee equations collapses the 4-port to the two-port
    seen from ports 1 and 2:
    \[ [S]_{\text{shorted}} = \begin{bmatrix} -1 & 0 \\ 0 & -1 \end{bmatrix} \]
\end{itemize}}{%
\begin{itemize}
  \item \textbf{Behaviour}: each collinear arm reflects \textbf{all} of its incident power
    with a $180^\circ$ phase reversal, and the two arms stay \textbf{completely isolated}
    from each other.
  \item The E-arm and H-arm reflections add in the driven arm and cancel in the other one,
    which is why $S_{21}$ stays zero.
  \item Still lossless (each column sums to 1), still reciprocal. The tee has become a
    \textbf{pair of independent short circuits}, useful as a reflective phase shifter or a
    reference short.
\end{itemize}}

\penalty0\vspace{2pt}

%=======================================================================
\T{3.7 Identifying a Device from its S-Matrix}

\Q{Identify and explain, with a neat diagram, the passive microwave device described by the given S-matrix}{\m{5}\m{6}\m{8}\m{10}}
{\color{sub}\footnotesize \tF{5} \yr{\textbf{74 Bh}, \textbf{75 Bh}, \textbf{78 Ch}, \textbf{\texttt{81 Bh}}, \textbf{\texttt{82 Bh}}} \gap$\cdot$\gap \mk{in all five papers the printed matrix is a magic tee}}

\sbs{%
\lead{The six-step method}
\begin{enumerate}[label=\arabic*., leftmargin=1.4em]
  \item \textbf{Count the ports}: the order of the matrix is the port count. $4\times4$
    narrows it to coupler, magic tee or hybrid ring immediately.
  \item \textbf{Read the diagonal}: zeros mark the matched ports. An all-zero diagonal on
    a 4-port means matched everywhere.
  \item \textbf{Test the transpose}: symmetric means passive and isotropic; asymmetric
    means ferrite or active.
  \item \textbf{Test the columns}: sums of 1 and orthogonal pairs mean lossless. A
    shortfall is dissipation.
  \item \textbf{Map the zeros}: an off-diagonal zero is an \textbf{isolated pair}. Which
    pairs are decoupled is the strongest single clue.
  \item \textbf{Read the signs and the $j$'s}: a sign flip down a column means an
    \textbf{antiphase (difference, E) arm}; matching signs mean an \textbf{in-phase (sum,
    H) arm}; a factor $j$ means a \textbf{quadrature} coupled arm.
\end{enumerate}
\lead{Answer skeleton}
\begin{itemize}
  \item Name the device in the first line. Do not narrate the tests first.
  \item Draw it: the 3D waveguide sketch with all arms labelled, plus the schematic.
  \item Justify with the four tests, one line each.
  \item Close with two uses and one limitation (the tee needs matching irises at the
    collinear arms to reach the ideal matrix in practice).
\end{itemize}}{%
\lead{Fingerprint table}
\begin{tabularx}{\linewidth}{@{}L{3.05cm}X@{}}
\toprule
\hrow \thead{Pattern} & \thead{Device} \\
\midrule
$2\times2$, $S_{11}=S_{22}=0$, $S_{12}=S_{21}=1$ & Matched lossless line or ideal matched network \\
$2\times2$, $S_{12}=0$, $S_{21}=1$ & Isolator \\
$2\times2$, $|S_{21}|>1$, $|S_{12}|$ tiny & Amplifier (active, not reciprocal) \\
$3\times3$, cyclic ones, not symmetric & Circulator \\
$3\times3$ with a $\pm1/\sqrt{2}$ column & E-plane tee ($-$) or H-plane tee ($+$) \\
$4\times4$, zero diagonal, two zero $2\times2$ blocks, one column sign-flipped & \mk{Magic tee} \\
$4\times4$, zero diagonal, $j$ on the coupled arm, one zero per row & Directional coupler \\
\bottomrule
\end{tabularx}}

\sbs{%
\lead{The four matrices the papers actually print}
\begin{itemize}
  \item \textbf{78 Ch} \m{10}, printed without its normalising factor:
    \[ [S] = \begin{bmatrix} 0&0&1&1 \\ 0&0&-1&1 \\ 1&-1&0&0 \\ 1&1&0&0 \end{bmatrix} \]
    Zero diagonal, symmetric, the $(1,2)$ and $(3,4)$ blocks zero, column 3 sign-flipped:
    \textbf{magic tee}. It is unitary only after dividing by $\sqrt{2}$, and saying so
    earns the ``judge the condition'' mark.
  \item \textbf{\texttt{81 Bh}} \m{5} and \textbf{74 Bh} \m{8}, printed symbolically:
    \[ [S] = \begin{bmatrix}
       S_{11}&S_{12}&S_{13}&S_{14} \\ S_{12}&S_{22}&S_{13}&-S_{14} \\
       S_{13}&S_{13}&0&0 \\ S_{14}&-S_{14}&0&0 \end{bmatrix} \]
    Column 3 in phase (H-arm), column 4 antiphase (E-arm), $S_{33}=S_{44}=0$:
    the \textbf{magic tee before the unitary property is applied}. Finish the job and
    $S_{11}=S_{22}=S_{12}=0$, $S_{13}=S_{14}=1/\sqrt{2}$.
\end{itemize}}{%
\begin{itemize}
  \item \textbf{75 Bh} \m{6}, the same device with the collinear isolation already
    written in:
    \[ [S] = \begin{bmatrix}
       S_{11}&0&S_{13}&S_{14} \\ 0&S_{22}&-S_{13}&S_{14} \\
       S_{13}&-S_{13}&0&0 \\ S_{14}&S_{14}&0&0 \end{bmatrix} \]
    Here column 3 is the antiphase E-arm and column 4 the in-phase H-arm, the mirror of
    the previous labelling. \mk{Read the signs, not the port numbers.}
  \item \textbf{\texttt{82 Bh}} \m{4} prints two \emph{amplifier} sets instead, so steps
    3 and 4 answer it: $S_{12} \neq S_{21}$ and $|S_{21}| > 1$ make both of them active
    and neither reciprocal ($\S$3.4, Problem 2).
\end{itemize}}

\penalty0\vspace{2pt}

%=======================================================================
\T{3.8 Power Splitting and Combining: the Radar Questions}

\Q{Two identical radar transmitters must deliver twice the power to one antenna; and an ASR needs half the transmitter power at the antenna through a duplexer. Solve both with an S-matrix}{\m{6}\m{8}\m{10}}
{\color{sub}\footnotesize \tF{4} \yr{73 Ma, \textbf{73 Bh}, 74 Ma, \textbf{76 Bh}} \gap$\cdot$\gap \mk{one device, two directions: the magic tee combines and splits}}

\sbs{%
\lead{Combining: two transmitters into one antenna}
\begin{itemize}
  \item \textbf{Device}: magic tee (or a $3$ dB hybrid coupler).
  \item \textbf{Connection}: the two transmitters drive the collinear arms 1 and 2
    \textbf{in phase and at equal amplitude}; the antenna hangs on the H-arm (port 4);
    the E-arm (port 3) carries a \textbf{matched dummy load}.
  \item \textbf{Mathematics}, with $a_1 = a_2 = a$ and $a_3 = a_4 = 0$:
    \[ b_4 = \tfrac{1}{\sqrt{2}}(a_1 + a_2) = \sqrt{2}\,a
       \;\Rightarrow\; |b_4|^2 = 2|a|^2 \]
    \[ b_3 = \tfrac{1}{\sqrt{2}}(a_1 - a_2) = 0 \]
  \item \textbf{Result}: the antenna receives $2P$, twice what either transmitter can
    deliver, and the dummy load receives nothing.
  \item \textbf{Why the tee and not a plain T-junction}: $S_{12} = 0$, so the two
    transmitters are \textbf{isolated} and cannot pull or damage each other; and all four
    ports stay matched, so neither transmitter sees a VSWR.
  \item \textbf{Practical riders} worth a mark each:
  \begin{itemize}
    \item Phase must be equalised, so the two feed lines are cut to the same electrical
      length.
    \item Amplitude imbalance $\Delta$ shows up as $|b_3|^2 \neq 0$: the dummy load
      absorbs exactly the mismatch, which is what protects the transmitters.
    \item The dummy load must be rated for the full unbalanced power.
  \end{itemize}
\end{itemize}}{%
\lead{Splitting: half the transmitter power to the antenna}
\begin{itemize}
  \item \textbf{Device}: the same magic tee, used as a duplexer, driven from the H-arm.
  \item \textbf{Connection}: transmitter on port 4 (H-arm), antenna on port 1, matched
    load on port 2, receiver on port 3 (E-arm).
  \item \textbf{Mathematics}, with $a_4 = a$ and all other inputs zero:
    \[ b_1 = \tfrac{1}{\sqrt{2}}a, \quad b_2 = \tfrac{1}{\sqrt{2}}a
       \;\Rightarrow\; |b_1|^2 = |b_2|^2 = \tfrac{1}{2}|a|^2 \]
    \[ b_3 = 0 \]
  \item \textbf{Result}: the antenna gets exactly $P/2$, as \textbf{76 Bh} requires, and
    the receiver on the E-arm gets \textbf{none} of the transmitter power, which is the
    whole point of a duplexer.
  \item \textbf{On receive}, the echo arrives at port 1 alone and divides between the
    E-arm and the H-arm, so half of it reaches the receiver. The $3$ dB penalty in each
    direction is the price of the magic-tee duplexer, and naming it shows you understand
    the device rather than reciting it.
  \item \textbf{Better answer for a real ASR}: a \textbf{ferrite circulator} duplexer
    routes the full power $1 \rightarrow 2$ and the full echo $2 \rightarrow 3$ with no
    $3$ dB loss. Quote it as the improvement whenever the question does not force the
    half-power split.
\end{itemize}
\lead{Answer skeleton for both}
\begin{itemize}
  \item Name the device, draw it with every arm labelled and every termination named.
  \item Write $[S]$ for the magic tee.
  \item Write the drive vector $[a]$ the connection creates.
  \item Multiply out, and read the answer off $|b|^2$.
  \item Close with the isolation and matching argument: that is the ``sufficient
    explanation'' the marks are for.
\end{itemize}}

\penalty0\vspace{2pt}

%=======================================================================
\T{3.9 PYQ Mapping}

\lead{Theory questions, covered in this document}
\begin{itemize}
  \item \textbf{Why S-parameters; importance in microwave network analysis; vs
    h-parameters} $\rightarrow$ $\S$3.1.
    \mk{The most repeated topic in the chapter, 13 papers.}
    \yr{69 Bh \m{4}, \textbf{70 Bh} \m{4}, 70 Ma \m{4}, \textbf{80 Ch} \m{5},
    \textbf{\texttt{81 Bh}} \m{3}}
  \item \textbf{Derive S-parameters for a two-port network} $\rightarrow$ $\S$3.2.
    \yr{69 Bh \m{6}, \textbf{70 Bh} \m{4}, \texttt{80 Ba} \m{6}, \textbf{80 Ch} \m{5}}
  \item \textbf{Properties of the S-matrix} $\rightarrow$ $\S$3.3. The four tests are
    also the identification tool of $\S$3.7. \yr{\textbf{\texttt{80 Bh}} \m{4}}
  \item \textbf{Return loss and insertion loss} $\rightarrow$ $\S$3.4.
    \yr{\textbf{72 Ash} \m{2}, \textbf{\texttt{79 Bh}} \m{1}}
  \item \textbf{Three-port network: analysis, properties, characteristic parameters}
    $\rightarrow$ $\S$3.5, with the impossibility theorem and its three escapes.
    \yr{\textbf{72 Ash} \m{5}, 73 Ma \m{4}, \textbf{73 Bh} \m{4}, 74 Ma \m{6},
    \textbf{\texttt{82 Bh}} \m{4}}
  \item \textbf{Three-port parameters using an E-plane tee example} $\rightarrow$ $\S$3.5
    for the parameters, $\S$3.6 for the matrix, companion Problem 5 for the derivation.
    \yr{\textbf{\texttt{80 Bh}} \m{4}}
  \item \textbf{S-matrices of E-plane tee, H-plane tee and a duplexer, and their uses}
    $\rightarrow$ $\S$3.6. \yr{\texttt{81 Ba} \m{8}}
  \item \textbf{Magic tee with both E- and H-arms shorted} $\rightarrow$ $\S$3.6,
    \added{} derived here because no source note carries it.
    \yr{\textbf{\texttt{82 Bh}} \m{4}}
  \item \textbf{Identify the passive device from a printed S-matrix} $\rightarrow$
    $\S$3.7, with all four printed matrices worked.
    \yr{\textbf{74 Bh} \m{8}, \textbf{75 Bh} \m{6}, \textbf{78 Ch} \m{10},
    \textbf{\texttt{81 Bh}} \m{5}}
  \item \textbf{Radar power combining and duplexer power splitting} $\rightarrow$ $\S$3.8.
    \yr{73 Ma \m{8}, \textbf{73 Bh} \m{8}, 74 Ma \m{6}, \textbf{76 Bh} \m{10}}
  \item \textbf{Evaluate the S-matrix of an amplitude modulator} $\rightarrow$ $\S$3.6 and
    companion Problem 6. \added{} \yr{\textbf{78 Ch} \m{7}}
\end{itemize}

\lead{Numerical content}
\begin{itemize}
  \item \textbf{This chapter has a numerical companion}, \textit{Ch3 - RF and MW Network
    Theory and Analysis - Numerical Problems and Solutions}. It carries every S-parameter
    calculation the archive asks for: the \textbf{\texttt{79 Bh}} reciprocal / lossless /
    return-loss / reflected-power question, the \textbf{\texttt{82 Bh}} comparison of two
    amplifier sets, the four device-identification matrices, the E-plane, H-plane and
    magic-tee derivations, and the two radar power problems.
  \item \textbf{Not in this chapter}: amplifier stability ($K$, $\Delta$, $\mu$) and the
    unilateral / bilateral maximum gain calculations. Those ride on S-parameters but the
    archive files them under \textbf{RF Design Practices}, so they belong to Chapter 6.
    $\S$3.4 gives only the reading of the numbers that \textbf{\texttt{82 Bh}} asks for.
\end{itemize}

\lead{Syllabus coverage check}
\begin{itemize}
  \item \textit{Scattering matrix and its properties} $\rightarrow$ $\S$3.3, $\S$3.5,
    $\S$3.6 \checkmark
  \item \textit{S-parameter derivation and analysis} $\rightarrow$ $\S$3.1, $\S$3.2,
    $\S$3.4, $\S$3.7, $\S$3.8 \checkmark
\end{itemize}

\lead{Cross-references, so nothing is said twice}
\begin{itemize}
  \item \textbf{Chapter 1} $\S$1.2 justifies S-parameter analysis as part of the
    lumped-vs-distributed answer. The full case is $\S$3.1 here.
  \item \textbf{Chapter 2} $\S$2.8 carries the two-line S-matrix rider tacked onto ten
    stub-tuning questions ($[\Gamma_L]$ before matching, $\left[\begin{smallmatrix}0&1\\1&0\end{smallmatrix}\right]$
    after). The theory behind it is $\S$3.2 and $\S$3.3.
  \item \textbf{Chapter 4} covers the physical construction of the tees, couplers,
    circulators and waveguide junctions whose matrices are catalogued in $\S$3.6.
\end{itemize}

\penalty0\vspace{2pt}
